\chapter{huidige staat van documentatie}
\label{ch:huidige_staat_documentatie}

Alvorens dieper in te gaan op de tekortkomingen van de huidige aanpak, is het noodzakelijk om de huidige staat van documentatie binnen het EMAV-web project te schetsen en deze te contrasteren met de situatie in de originele Windows Forms applicatie. Momenteel bevindt het project zich in een overgangsfase, waarbij de lessen uit het verleden de koers voor de toekomst bepalen.

\section{Lessen uit het Legacy project}
In de oorspronkelijke EMAV-applicatie was de documentatie extreem gefragmenteerd. Voor de onderzoekers bestond de documentatie voornamelijk uit lijvige Word-documenten en PDF-handleidingen die losstonden van de codebasis. Voor de programmeurs waren er pogingen ondernomen om complexe business-flows (zoals de \textit{kunstmest-flow}) te documenteren in afzonderlijke Markdown-bestanden binnen de \texttt{docs}-map van de repository.

Deze aanpak leed echter aan een chronisch gebrek aan onderhoudbaarheid. In de praktijk fungeerden deze Markdown-bestanden als een soort "schaduw-code": ze bevatten grote blokken gekopieerde C\#-logica met minimale tekstuele uitleg. Dit zorgde voor een enorme \textit{documentation drift}: zodra een formule in de rekenkern werd aangepast, werd de corresponderende uitleg in de documentatiemap vergeten. Het resultaat was een systeem waarbij de documentatie eerder voor verwarring zorgde dan voor verduidelijking, wat de afhankelijkheid van de oorspronkelijke ontwikkelaars alleen maar vergrootte.

\section{Huidige aanpak in EMAV-web: Technische README's}
In de nieuwe EMAV-web applicatie is er een bewuste verschuiving merkbaar naar "Documentation as Code". De primaire vorm van documentatie voor de ontwikkelaars bestaat momenteel uit Markdown-bestanden, zoals de centrale \texttt{README.md} in de root van de repository. Deze bestanden zijn veel directer gekoppeld aan het ontwikkelproces: ze bevatten "TODO NEXT"-lijsten, instructies voor database-migraties en architecturale diagrammen in Mermaid-formaat.

Hoewel dit een aanzienlijke verbetering is voor de IT-onboarding, blijft de kloof met de onderzoekers van ILVO gapen. De informatie in deze bestanden is nagenoeg uitsluitend technisch van aard. Voor een onderzoeker die wil begrijpen hoe een emissiefactor tot stand komt, biedt een README met instructies over \texttt{WebDbContext} migraties geen enkel inzicht. De documentatie is "ontwikkelaar-tot-ontwikkelaar" geschreven, waardoor de wetenschappelijke essentie van het project nog steeds onderbelicht blijft.

\section{XML-commentaren en de limieten van DocFX}
Binnen de nieuwe C\#-codebase wordt consequent gebruikgemaakt van \texttt{/// <summary>} XML-commentaren. Er zijn concrete plannen om \textbf{DocFX} te gebruiken om deze commentaren automatisch te publiceren als een statische website. Hoewel dit de technische transparantie verhoogt, lost het de fundamentele vraag van de onderzoeker niet op. XML-commentaren zijn gebonden aan de structuur van de klassen en methoden. Ze zijn uitstekend voor het beschrijven van input-parameters, maar ongeschikt voor het weergeven van de complexe, domeinoverschrijdende wetenschappelijke redeneringen die EMAV uniek maken.

\section{De \textquotedbl Single Source of Truth\textquotedbl \@ ontbreekt}
De huidige realiteit is dat de eigenlijke "levende" documentatie van de wetenschappelijke modellen zich nog steeds buiten de broncode bevindt, in Excel-sheets en verslagen van commissies. Er is geen enkel mechanisme dat garandeert dat de formule in de Excel-sheet van de onderzoeker exact dezelfde is als de implementatie in \texttt{ILVO.EMAV.Core}. Deze versnippering is het kernprobleem dat in dit onderzoek wordt geadresseerd. De noodzaak voor een vorm van documentatie die zowel executabel is voor de machine als leesbaar voor de wetenschapper — zoals bij literate programming — komt direct voort uit deze historische en actuele tekortkomingen.
